\scanpage{88}
\deleted{et par suite $\widetilde\rho$} Considérons $R_0=\widetilde\rho_0\longrightarrow R$ et les flèches qui deviennent des isom, soit $\Sigma$, \deleted{\unsure{On a donc} $\widetilde\rho\Sigma^{-1}$} Soit \deleted{\unsure{$\rho_0$}} $\rho_1=R_0\Sigma^{-1}$, \inserted{$R_1=\widetilde{\rho_1}$} Alors $C\mapsto\underline{\mathrm{Hom}}_{\mathrm{Cat}}(R_0\Sigma^{-1},C)\simeq\underline{\mathrm{Hom}}_{\Delta}(R_1,C)$ est alg. $R_1\to R$. On construit une suite transfinie de catégories $R_\alpha\in\Delta$, et de $R_\alpha\to R$ ; toutes les \inserted{théories} $\sigma(R_\alpha)$ sont \unsure{algébrique|algébriques} (le passage aux \deleted{\unsure{limites|lim}}\inserted{ordinaux limites} se fait par le critère d), ou $d'$\deleted{\unsure{)|!}} , au choix). On suppose donc

\textcircled{$4^\circ$} Soit $R\in\mathrm{Ob}\,\Delta$, et $\deleted{\unsure{J}}I\xrightarrow{u}R$ un foncteur \deleted{\unsure{fidèle}}\inserted{(dans Cat.)} tel que $\forall\,C\in\mathrm{Ob}\,\Delta$, le foncteur $\underline{\mathrm{Hom}}_\Delta(R,C)\mathbin{\unsure{-|\to}}\underline{\mathrm{Hom}}_{\Delta}(R',C)$ fidèle ($u$ est \deleted{dit}\inserted{foncteur $\Delta$-}générateur). Construisons \deleted{un}\inserted{deux 2-}syst. inductif transfini de catégories $R_\alpha\in\mathrm{Ob}\,\Delta$, \inserted{$\rho_\alpha\in\mathrm{Ob}(\mathrm{Cat})$} tel que les morphismes de transition $R_\alpha\mathbin{-}R_\beta$ sont des $\Delta$-foncteurs, et \deleted{\unsure{un morphisme|les morphismes} \unsure{fonct|foncteurs}}liés en \unsure{a|à} 2-syst. inductif dans $R$ par des $u_\alpha:R_\alpha\longrightarrow R$ \inserted{$v_\alpha:\rho_\alpha\to R$} ($u_\alpha\in\mathrm{Fl}(\Delta)$), en posant
% slanted
\margin{faut il supposer $u$ pl fidèle ??}
\[
\begin{aligned}
R_0 &=(\widetilde I,\ \deleted{\unsure{\xi}}\ u_0:\widetilde I\mathbin{-}R\text{ prolongeant }u),\ \rho_0=R_0\\
R_{\alpha+1}&\ \inserted{\mathbin{-}R}\ \text{déduit de }R_\alpha\mathbin{-}R\text{ en définissant }
\Sigma_\alpha=\text{ens des}\\
&\qquad\text{flèches \inserted{de $R_\alpha$} qui deviennent inv. dans }R.\\
R_{\alpha+1}&=\underbrace{R_\alpha\Sigma^{-1}}_{\rho_{\alpha+1}}\quad\text{(au sens de }\Delta)
\end{aligned}
\]
% slanted
\margin{\deleted{prendre} \inserted{sous-}plutôt \inserted{partout ?} $R_{\alpha+1}=$ image essentielle de $R_\alpha\Sigma^{-1}$ dans $R$ (sous-cat. pleine de $R$) ?\par $\rho_{\alpha+1}\mathbin{-}R_{\alpha+1}\mathbin{-}\rho_{\alpha+2}\mathbin{-}R_{\alpha+2}$}
\[
 R_{\varinjlim\alpha_i}=\varinjlim_{\Delta}R_{\alpha_i}.
\]
Alors pour $\alpha$ assez grand \inserted{(dans l'univers $\mathcal U$)}, on a $R_\alpha\xrightarrow{\simeq}R$.

\uline{NB} Pour un ordinal limite, on voit que $\rho_\alpha\mathbin{-}R$ est conservatif \unsure{!!|?}. Si les $\Delta$-foncteurs \deleted{envisagés sont par} \deleted{\unsure{soit un}} commutent aux noyaux (\unsure{E|et|\ensuremath{=}} \unsure{que x|ceux|quot.} \unsure{et|ci} existent) ou conoyaux (itou), cel implique que les foncteurs $R_\alpha$ sont \deleted{pl. fidèl} \uline{fidèles}, donc les $R_\alpha$ sont des sous-catégories. Pour tt cardinal $c$, si $\alpha$ est plus élevé que $c$, $R_\alpha$ est stable sous les limites de types $\delta'\in d'$ ou $\delta''\in d'$ pour $\mathrm{card}(\delta'),\mathrm{card}(\delta'')\leq c$. Donc si $d'$ et $d''$ sont \unsure{petits|\deleted{petits}}
% slanted
\margin{avec le choix \unsure{primo|prim.} proposé on voit qu'on aura \deleted{des} \inserted{pour $\alpha$ \unsure{ind.} lim.} foncteurs pl fidèles $\rho_\alpha\subset R$, la question est si $R$ est égal à $\rho_\alpha$ pour $\alpha$ assez grand\ldots}
% archive sheet number: 87
